\documentclass[10pt,a4paper]{amsart}
\usepackage[margin=19mm]{geometry}
\usepackage{amsmath,amssymb,mathtools}
\usepackage[scr=rsfs,cal=euler]{mathalfa}
\usepackage{xfrac}
\usepackage[unicode,hidelinks,bookmarks=false]{hyperref}
\newcommand{\ie}{i.e.\ }
\newcommand{\cf}{cf.\ }
\newcommand{\resp}{respectively\ }
\newcommand{\Subsection}{\subsection}
\providecommand{\nobreakdash}{\nobreak\hbox{-}}
\setlength{\emergencystretch}{2em}
\multlinegap0pt
\usepackage{bm}
\usepackage{bbm}
\usepackage{dsfont}
\usepackage{xspace}
\usepackage{cancel}
\usepackage{mathtools}
\usepackage{amsthm}
\makeatletter
\@ifundefined{definition}{\newtheorem{definition}{Definition}}{}
\@ifundefined{theorem}{\newtheorem{theorem}{Theorem}}{}
\@ifundefined{corollary}{\newtheorem{corollary}[theorem]{Corollary}}{}
\makeatother
\newcommand{\E}{\nu}
\newcommand{\Green}{\ensuremath{\mathbf{G}}}
\newcommand{\PP}{\nu}
\newcommand{\RR}{\mathbb{R}}
\newcommand{\exact}{\star}
\renewcommand{\c}[1]{\ensuremath{\mathcal{#1}}}
\title[On the duality between height functions and continuous spin ]{On the duality between height functions and continuous spin models (Theorem 3.2)}
\author{Diederik van Engelenburg, Marcin Lis}
\date{Source: arXiv:2303.08596v1 (CC BY 4.0)}
\begin{document}
\maketitle

\noindent\emph{Source and licence.} On the duality between height functions and continuous spin models. Version arXiv:2303.08596v1 (CC BY 4.0), distributed under the Creative Commons Attribution 4.0 International licence, as recorded in the arXiv licence metadata for this exact version and shown on the abstract page https://arxiv.org/abs/2303.08596v1. Full rights notice: licenses/arxiv-2303.08596-CC-BY-4.0.txt.\par\smallskip

\noindent\textit{Editorial scope.} Counted unit: the Theorem printed as Theorem 3.2 of the article (source0.tex lines 522--526, source label \texttt{thm:extransinv}), delivered together with the article's own complete original proof of it (source0.tex lines 528--563, one \texttt{proof} environment of 36 lines). No mathematics is written, repaired or re-derived: statement and proof are the article's own text line for line. Required context retained uncounted: def:potential (lines 275--286); cor:upper (lines 450--456). Every definition, display and label of those lines is retained unchanged. Omitted from this package and not redistributed: 1--274, 287--449, 457--521, 527--527, 564--1555. Nothing inside a retained excerpt is omitted or altered: every retained body is delivered exactly as the article's lines are written, so no omission receipt of any kind accompanies this package. Citations: the retained excerpts carry no citation command at all, so no bibliography entry of the article is transcribed and no reference is redistributed. Reference adaptation: none was needed and none was made; every reference inside the delivered unit resolves inside it, so no unresolved reference or citation is produced. Printed slips kept literal: no printed word, symbol, hypothesis or number of the counted statement or of its proof was altered. Layout and class: the article's own class is \texttt{amsart} with packages including geometry, graphicx, subcaption, stmaryrd, amssymb, amsthm, dsfont, hyperref, mathrsfs, amsrefs, amsmath, centernot, enumerate, xcolor; the wrapper is the lane's pinned \texttt{amsart} editorial wrapper at 10pt on A4, which restates the theorem-like environments the retained text needs and adds the article's own macro definitions that the retained text needs (\texttt{\textbackslash E}, \texttt{\textbackslash Green}, \texttt{\textbackslash PP}, \texttt{\textbackslash RR}, \texttt{\textbackslash c}, \texttt{\textbackslash exact}), with extra wrapper packages (bm, bbm, dsfont, xspace, cancel, mathtools). The delivered numbering is the unit's own. Assets: no retained span contains a figure environment, a graphics-inclusion command, a PDF-page inclusion, a table or a TikZ picture; no image file is copied into this package, the package declares no asset dependency and no image file is redistributed with it. Licence: version arXiv:2303.08596v1 of this article is distributed under the Creative Commons Attribution 4.0 International licence, as recorded in the arXiv licence metadata for this exact version and shown on the abstract page https://arxiv.org/abs/2303.08596v1. A full rights notice is delivered at licenses/arxiv-2303.08596-CC-BY-4.0.txt. Characters: every retained excerpt is pure ASCII, so the delivered engine, XeLaTeX with its default Latin Modern text font, reports no missing character.
\begin{definition}[Height function and spin potentials] \label{def:potential}
Let $\mathcal V: \mathbb Z\to \mathbb R \cup \{ +\infty\}$ be symmetric, i.e., $\mathcal V (n)=\mathcal V(-n)$, such that 
\begin{align} \label{eq:Vsum}
\sum_{n\in \mathbb Z}n^2 \exp(-\mathcal V(n))<\infty,
\end{align}
and moreover such that $\exp(-\mathcal V)$
is \emph{positive definite}: for all $\alpha\in \mathbb R$, 
\begin{align}\label{eq:posdef}
	w(\alpha):=\exp(-\mathcal V(0))+\sum_{n=1}^{\infty}\exp(-\mathcal V(n))2\cos( n\alpha) > 0.
\end{align}
We call $\mathcal V$ the \emph{height function potential}, and $\mathcal U (\alpha) := -\log w(\alpha)$ the \emph{spin potential}.
\end{definition}

\begin{corollary}[GFF upper bound on variance]\label{cor:upper}
	For any $f\in \Omega^0_0(\mathbb R)$,
	\begin{align*}
		\nu_{\exact}[(h, f)^2] \leq  C  ( \Delta^{-1}f,f)_{\Omega^0}=C\sum_{v,v'\in V} f_vf_{v'}\frac{\Green(v, v')}{\deg(v')},
	\end{align*}
	where $C$ is as above.
\end{corollary}

\begin{theorem} \label{thm:extransinv}
	Let $\Gamma = (\mathscr V,\mathscr E)$ be a transient graph and $\c{V}$ a height function potential as in Definition~\ref{def:potential}. 
	Then, there exists an infinite volume Gibbs measure on $\Gamma$ with respect to $\mathcal V$. If $\Gamma$ is moreover an amenable Cayley graph, 
	there exist translation invariant Gibbs measures. 
\end{theorem}

\begin{proof}
Let now $G_N\nearrow \Gamma$, as $N\to \infty$ be an exhaustion of $\Gamma$ by finite subgraphs $G_N$. 
Define the boundary $\partial_N := \partial G_N$ to be the set of vertices in $ G_N$ adjacent to a vertex from outside of $ G_N$.
Let $\E_{G_N, \exact}[\cdot]$ be the expectation associated with the height function on $ G_N$ with $0$-boundary conditions. 
Fix any $\Lambda \subset \c{V}$ finite and let $f: \Lambda \to \RR$ be any function.
It follows from Corollary~\ref{cor:upper} that 
\[
	\E_{G_N,\exact}[(h, f)^2] \leq C \max_{v \in \Lambda} \frac{\Green_{N} (v,v)}{\deg(v)} (f, f)_{\Lambda}. 
\]
where $C < \infty$, and $\Green_N$ is the Green's function of simple random walk on $G_N$ killed on hitting $\partial_N$. 
Since $\Gamma$ is transient, the right-hand side is uniformly bounded in $N$. 
Therefore, the sequence $\PP_{G_N, \exact}(h |_{\Lambda})$ is tight and subsequential limits exist by Prokhorov's theorem. 
By a diagonal argument, we can extract a further sub-sequence $N_K$ so that $\PP_{G_{N_K}, \exact}(h|_{\Lambda})$ converges for each
$\Lambda$ finite. 
Any such subsequential limit is a Gibbs measure as it satisfies the DLR relations. 
This proves the first part of the theorem. 

For the second part, suppose that $\Gamma$ is an amenable Cayley graph, 
so that 
\(
	\E_{G_N, \exact}[h_v^2] \leq C',
\)
for some $C' < \infty$ which is independent of $v$ and the exhaustion $(G_N)_{N\geq 1}$. 
Let $\mu$ be a subsequential limit (which exists by the argument above, and is a Gibbs measure). 
Let $o \in \mathscr V$. Since $\Gamma$ is amenable, there is some F\o lner sequence (also called Van Hove sequence) $(F_N)_{N \geq 1}$ {of sets of vertices} containing $o$. 
This means $F_N \nearrow \mathscr V$ and $|\partial F_N| / |F_N| \to 0$ as $N \to \infty$. 
Let
\[
	\nu_N := \frac{1}{|F_N|} \sum_{x \in F_N} \mu \circ \theta_x, 
\]
where $\theta_x$ is the shift towards $x$ (since $\Gamma$ is a Cayley graph, this is the same as left multiplication in the group).
This is a Gibbs measure because the set of Gibbs measures is closed under translations and convex combinations. 
Moreover, $\nu_N[h_v^2] \leq C'$ for each $N$ and $v \in \mathscr V$. Therefore, $(\nu_N)_{N\geq 1}$ is tight. 
Let $\nu$ be any subsequential limit, which is again a Gibbs measure. 
By construction and since $|\partial F_N| / |F_N| \to 0$, we have $\nu \circ \theta_x = \nu$ for each $x$, and hence $\nu$ is translation invariant.
\end{proof}

\end{document}
